\documentclass{article}
\usepackage{graphicx}

\usepackage[utf8]{inputenc}
\usepackage[spanish]{babel}

\title{Actividad Complementaria - Display de 7 Segmentos}
\author[1]{Victor Eduardo Cano Mireles}
\author[1]{Ian Gerardo Esquivel Garcia}
\affil[1]{Departamento de Ingeniería Eléctrica y Computación, Universidad Autónoma de Ciudad Juárez}
\date{28 de abril de 2026}

\begin{document}

\maketitle  


\newpage
\section{Instrucciones}
    \begin{enumerate}
    \item{Identificar los requisitos del sistema mediante una lista:}
    \begin{itemize}
        \item{¿Qué debe hacer el sistema?}
        \item{¿Qué entradas tendrá?}
        \item{¿Qué salidas genera?}
    \end{itemize}
    \item{Definir el léxico del sistema con al menos 5 términos}
    \item{Elaborar el modelo entrada-proceso-salida}
    \item{Realizar la tabla de verdad del display}
    \item{Conectar el display de 7 segmentos al Arduino}
    \item{Programar Arduino para mostrar números del 0 al 9}
    \item{Crear una página web que contenga:}
    \begin{itemize}
        \item{Botones del 0 al 9}
        \item{Visualización del número seleccionado}
    \end{itemize}
    \item{Programar un servidor que:}
    \begin{itemize}
        \item{Reciba el número desde la web}
        \item{Envíe el dato por comunicación serial (9600 baudios)}
    \end{itemize}
    \item{Verificar que:}
    \begin{itemize}
        \item{El número cambia en la página web}
        \item{El número cambia en el display físico}
    \end{itemize}
    \item{Elaborar los siguientes modelos del sistema:}
    \begin{itemize}
        \item{Diagrama de estados}
        \item{Caso de uso}
        \item{Diagrama de actividades}
    \end{itemize}
    \item{Redactar una historia de usuario}
    \item{Describir el prototipo del sistema (hardware y software).}
    \end{enumerate}
\section{Requisitos del sistema}
    \subsection{¿Qué debe hacer el sistema?}
    \subsection{¿Qué entradas tendrá?}
    \subsection{¿Qué salidas generá?}

\end{document}
